%! Linear Combinations of Vectors — Unit 1, Lesson 1.1 — Homework (Answer Key)
\documentclass[10pt]{article}
\usepackage{linalg-article}
\usepackage{linalg-key}   % pulls in -boxes; do NOT also load -boxes
\newcommand{\TallMath}[1]{$\displaystyle #1\rule[-1.4em]{0pt}{3.2em}$}
\newcommand{\vv}{\mathbf{v}}
\newcommand{\ww}{\mathbf{w}}

\begin{document}

\pageheader{Unit 1, Lesson 1.1}{Homework --- Answer Key}
\namedateperiod

\vspace{0.12in}

\begin{notesbox}{Practice}
Throughout, let $\vv = \begin{bmatrix} 1 \\ 2 \end{bmatrix}$ and
$\ww = \begin{bmatrix} 3 \\ 1 \end{bmatrix}$. Show your steps.

\begin{enumerate}[leftmargin=1.9em, itemsep=10pt]
  \item Compute each linear combination.
    \begin{enumerate}[label=(\alph*), leftmargin=1.8em, itemsep=4pt]
      \item $\vv + \ww = {}$\ans{$\begin{bsmallmatrix} 4 \\ 3 \end{bsmallmatrix}$} \hfill
      \item $2\vv - \ww = {}$\ans{$\begin{bsmallmatrix} -1 \\ 3 \end{bsmallmatrix}$} \hfill
      \item $3\vv + \ww = {}$\ans{$\begin{bsmallmatrix} 6 \\ 7 \end{bsmallmatrix}$} \hfill
      \item $-2\vv + 2\ww = {}$\ans{$\begin{bsmallmatrix} 4 \\ -2 \end{bsmallmatrix}$}
    \end{enumerate}
  \item Find weights $c,d$ so that $c\,\vv + d\,\ww$ equals the given target. Show the match-up.
    \begin{enumerate}[label=(\alph*), leftmargin=1.8em, itemsep=4pt]
      \item \ans{$c+3d=5,\ 2c+d=5 \Rightarrow c=2,\ d=1$} \hfill
      \item \ans{$c+3d=7,\ 2c+d=9 \Rightarrow c=4,\ d=1$}
    \end{enumerate}
  \item For each pair of vectors, does it span the \textbf{whole plane} or only a
        \textbf{line}? Give a one-line reason.
    \begin{enumerate}[label=(\alph*), leftmargin=1.8em, itemsep=4pt]
      \item $\begin{bsmallmatrix} 1 \\ 3 \end{bsmallmatrix},\ \begin{bsmallmatrix} 2 \\ 6 \end{bsmallmatrix}$ \hfill
      \item $\begin{bsmallmatrix} 2 \\ 1 \end{bsmallmatrix},\ \begin{bsmallmatrix} 1 \\ 3 \end{bsmallmatrix}$ \hfill
      \item $\begin{bsmallmatrix} 4 \\ -2 \end{bsmallmatrix},\ \begin{bsmallmatrix} -2 \\ 1 \end{bsmallmatrix}$
    \end{enumerate}
    \ansline{(a) Line --- $\begin{bsmallmatrix} 2\\6 \end{bsmallmatrix}=2\begin{bsmallmatrix} 1\\3 \end{bsmallmatrix}$ (parallel).
    \ (b) Plane --- not parallel.
    \ (c) Line --- $\begin{bsmallmatrix} -2\\1 \end{bsmallmatrix}=-\tfrac12\begin{bsmallmatrix} 4\\-2 \end{bsmallmatrix}$ (parallel).}
  \item A juice bar stores two base juices as vectors
        $\begin{bsmallmatrix} \text{sugar (g)} \\ \text{vitamin C (mg)} \end{bsmallmatrix}$:
        $\mathbf{a}=\begin{bmatrix} 8 \\ 30 \end{bmatrix}$ and
        $\mathbf{b}=\begin{bmatrix} 5 \\ 50 \end{bmatrix}$.
    \begin{enumerate}[label=(\alph*), leftmargin=1.8em, itemsep=5pt]
      \item A new drink is the blend $2\mathbf{a}+\mathbf{b}$. Find its sugar and vitamin~C.
            \ans{$2\begin{bsmallmatrix} 8\\30 \end{bsmallmatrix}+\begin{bsmallmatrix} 5\\50 \end{bsmallmatrix}
            = \begin{bsmallmatrix} 21\\110 \end{bsmallmatrix}$: $21$ g sugar, $110$ mg vitamin C}
      \item \emph{Interpret:} what does the weight $2$ on $\mathbf{a}$ mean for this recipe?
            \ansline{The blend uses two parts of juice $\mathbf{a}$ for every one part of $\mathbf{b}$.}
    \end{enumerate}
\end{enumerate}
\end{notesbox}

\begin{extensionbox}
Recall a vector can be a \textbf{data record}. A smoothie is stored as
$\begin{bsmallmatrix} \text{cups fruit} \\ \text{cups yogurt} \end{bsmallmatrix}$, with base
recipes $\mathbf{r}=\begin{bsmallmatrix} 2 \\ 1 \end{bsmallmatrix}$ and
$\mathbf{s}=\begin{bsmallmatrix} 1 \\ 2 \end{bsmallmatrix}$. Write the blend $3\mathbf{r}+2\mathbf{s}$,
and explain what a linear combination \emph{means} when the vectors are recipes instead of
arrows. \ansline{$3\mathbf{r}+2\mathbf{s}=\begin{bsmallmatrix} 8\\7 \end{bsmallmatrix}$: use $3$
batches of recipe $\mathbf{r}$ and $2$ of $\mathbf{s}$, giving $8$ cups fruit and $7$ cups
yogurt. The weights are ``how many batches of each recipe'' you pour in.}
\end{extensionbox}

\begin{spiralbox}
\textbf{Coming next --- Lesson 1.2: Lengths and Angles from Dot Products.} You now combine
vectors and know which points they can reach. Next we measure \emph{how much two vectors point
the same way} with a single number --- the dot product --- which also recovers a vector's
length as a special case.
\end{spiralbox}

\end{document}
